% ===========================================================================
\subsection{Two-piece hexagonal cavity: Ta/Nb coating only}\label{sec:tanb}
% ===========================================================================
%
% STATUS: measured, PPMS, zero field. Q_0 = 1.4e6 at 2 K on the two-piece
% hexagonal body, 500 nm DC Ta + 750 nm DC Nb, 200 C. Its own bare-Cu baseline
% is 55,000. Q(T) falls to 13,000 at 9 K. No Q vs B taken.
% Remaining \TODO{} items are photographs and two secondary details.
%
% ===========================================================================

Nb$_3$Sn is the harder material to make, and every defect in the Ta/Nb underlayer propagates into it. Before pushing further on Nb$_3$Sn we therefore coated the two-piece hexagonal copper cavity with the Ta/Nb stack alone and stopped there. Niobium is a single-element sputtered film. It needs no Cu-Sn evaporation, no \qty{715}{\degreeCelsius} reaction, and no APS etch. Its $T_c$ of \qty{9.2}{\kelvin} and its $H_{c2}$ below \qty{0.4}{\tesla} disqualify it from a \qty{9}{\tesla} haloscope, so this cavity is not a candidate detector. It measures how good the copper body, the seam, the surface preparation, and the deposition system can be once the film chemistry is taken out of the problem.

\subsubsection{Coating}

The stack was \qty{500}{\nano\meter} of Ta followed by \qty{750}{\nano\meter} of Nb, both by DC magnetron sputtering onto a substrate held at \qty{200}{\degreeCelsius}. Ta was deposited at \qty{60}{\watt} in \qty{10}{\milli\torr} Ar, Nb at \qty{250}{\watt} in \qty{8}{\milli\torr} Ar. The Nb thickness exceeds three penetration depths, so the film screens the copper beneath it and the measurement reports the film rather than the substrate.

\begin{figure}[htb]
    \centering
    \TODO{PHOTO: interior of the Ta/Nb-coated two-piece hexagonal cavity, both halves. Worth showing the endcap regions specifically, since this is the direct visual contrast with Fig.~\ref{fig:EdgeEffect2piece}.}
    \caption{Two-piece hexagonal copper cavity after Ta/Nb deposition. \TODO{Note coverage at the endcaps and at the steeply inclined edges, and compare with the Nb$_3$Sn cavity of Fig.~\ref{fig:EdgeEffect2piece}.}}
    \label{fig:tanbCavity}
\end{figure}

\subsubsection{Results}

The cavity was measured in the PPMS at zero applied field. At \qty{2}{\kelvin} the unloaded quality factor reached $Q_0 = 1.4\times10^{6}$, which for $G = \qty{350}{\ohm}$ corresponds to $R_s = \qty{250}{\micro\ohm}$. The same body measured bare gave $Q = 55{,}000$ at its best, so the coating improved this specific cavity by a factor of \num{25}. It also exceeds the simulated ceiling for a seamless RRR-300 copper cavity of this geometry ($Q = 76{,}000$) by a factor of \num{18}, and the fully corrected Nb$_3$Sn result of Sec.~\ref{sec:2piecenb3sn} ($Q_0 \approx 1.4\times10^{5}$) by a factor of \num{10}.

Above the transition the ordering inverts. At \qty{9}{\kelvin} the cavity gives $Q = 13{,}000$, which is a quarter of its own bare-copper value and corresponds to $R_s = \qty{27}{\milli\ohm}$ against \qty{6.4}{\milli\ohm} for copper. Normal-state sputtered Nb is far more resistive than annealed copper, and because the film is thicker than its normal-state skin depth at \qty{11}{\giga\hertz}, the RF current sees Nb and not the copper underneath. The full swing across the transition is therefore a factor of \num{108} in $Q$, from $13{,}000$ at \qty{9}{\kelvin} to $1.4\times10^{6}$ at \qty{2}{\kelvin}. A transition placed at the bulk Nb value confirms two things at once: the sputtered film is clean, and the Ta barrier is keeping copper out of it.

\begin{figure}[htb]
    \centering
    \TODO{FIGURE: $Q$ (or $R_s$) versus temperature, \qty{2}{}--\qty{15}{\kelvin}, PPMS, zero field. Overlay the bare Cu curve for this same body and the Nb$_3$Sn cavity curve so all three sit on one axis. Mark $Q = 13{,}000$ at \qty{9}{\kelvin} and $Q_0 = 1.4\times10^{6}$ at \qty{2}{\kelvin}.}
    \caption{Quality factor versus temperature for the Ta/Nb-coated two-piece cavity, measured in the PPMS at zero field, with the bare copper and Nb$_3$Sn cavities shown for comparison. \TODO{State the transition midpoint read from the curve.}}
    \label{fig:tanbQvsT}
\end{figure}

\TODO{Surface preparation history for this body: hand polish only, or the LLNL etch and \qty{250}{\degreeCelsius} anneal as well? The bare-Cu figure of $55{,}000$ matches the fully etched and annealed state, so state which.}

\TODO{Substrate rotation during Ta and Nb deposition, and whether the endcaps were masked.}

No $Q$ versus field data were taken. Niobium enters the mixed state below \qty{0.4}{\tesla}, so a field sweep would measure vortex dissipation in a material that cannot serve the application, and the cavity was reserved for the zero-field baseline instead.

\subsubsection{What this fixes in the argument}

A copper body of this geometry, with this seam, prepared this way and coated in this deposition system, supports $Q_0 = 1.4\times10^{6}$. The Nb$_3$Sn cavity built on the same design and the same Ta/Nb foundation delivered a corrected $1.4\times10^{5}$. Everything the two measurements share cancels in that ratio: the copper, the seam, the antenna coupling, the geometric factor, the VNA. An order of magnitude therefore separates what the hardware permits from what the Nb$_3$Sn process achieved, and none of it is attributable to the cavity or to the measurement chain. It belongs to the three steps that convert Ta/Nb into Nb$_3$Sn --- thermal evaporation of Cu-Sn onto a concave surface without rotation or heating, reaction at \qty{715}{\degreeCelsius}, and the APS etch that stripped both endcaps.
